\documentclass[sigconf,10pt]{acmart}

%% Packages
\usepackage{booktabs}
\usepackage{graphicx}
\usepackage{subcaption}
\usepackage{xcolor}
\usepackage{multirow}
\usepackage{amsmath}

%% Remove ACM copyright block for draft
\setcopyright{none}
\settopmatter{printacmref=false}
\renewcommand\footnotetextcopyrightpermission[1]{}
\pagestyle{plain}

\begin{document}

\title{Untangling GPU Architecture Design for LLM Inference:\\
A Systematic Design Space Exploration}

\author{Anonymous Authors}
\affiliation{\institution{Anonymous Institution}}

\begin{abstract}
Large Language Model (LLM) inference presents fundamentally different
architectural demands between its two phases---prefill (prompt processing)
and decode (autoregressive token generation). Through systematic profiling
of 61 workload configurations and design space exploration of 7,875
architecture variants using the LLMCompass simulation framework, we uncover
several insights that challenge conventional GPU design wisdom. Our key
findings include: (1)~systolic array dimension is 31.7$\times$ more
area-efficient than core count for prefill acceleration; (2)~decode
latency is determined by a single variable---HBM bandwidth---with all
other parameters collectively accounting for $<$1\% of variation;
(3)~multi-device scaling for decode is severely sub-linear ($\sim$15\%
efficiency at 8 devices) and can \emph{increase} latency beyond 4 devices;
(4)~prefill and decode optimization vectors are anti-correlated
($r = -0.29$) on the Pareto frontier, motivating heterogeneous
prefill-decode architectures. We distill these findings into six actionable
design guidelines for next-generation LLM accelerators and demonstrate that
a restructured area allocation can achieve equivalent performance at
35--45\% smaller die area.
\end{abstract}

\keywords{GPU Architecture, LLM Inference, Design Space Exploration,
Roofline Model, Hardware-Software Co-Design}

\maketitle

%% ===========================================================================
\section{Introduction}
\label{sec:intro}

The deployment of Large Language Models (LLMs) at scale has become one of
the dominant workloads in modern datacenters. Models such as GPT-3
(175B parameters)~\cite{brown2020language}, LLaMA (7B--70B)~\cite{touvron2023llama},
and their successors require massive compute and memory bandwidthfor inference serving. As LLM inference becomes a first-class design target
for GPU vendors, understanding the precise architectural bottlenecks and
tradeoffs is critical.

LLM inference consists of two distinct phases with fundamentally different
computational characteristics. The \textbf{prefill} phase processes the
input prompt through all transformer layers, performing large matrix
multiplications that are typically compute-bound. The \textbf{decode} phase
generates tokens one at a time, executing the same model architecture but
with much smaller tensor dimensions, resulting in memory-bandwidth-limited
execution.

Despite this well-known dichotomy, the quantitative architectural
implications remain underexplored. Specifically:
\begin{itemize}
\item \emph{Which architectural parameters matter for each phase, and by how much?}
\item \emph{What are the area-performance tradeoffs, and where do diminishing returns occur?}
\item \emph{How does multi-device scaling behave differently for each phase?}
\item \emph{Is it possible to design a single architecture that is optimal for both?}
\end{itemize}

In this work, we conduct a systematic study using LLMCompass~\cite{llmcompass},
a hardware-software co-simulation framework that models GPU execution at the
operator level using a roofline-based performance model and estimates die area
through a process-aware cost model. We make the following \textbf{contributions}:

\begin{enumerate}
\item \textbf{Comprehensive bottleneck taxonomy} across 61 workload
configurations spanning 5 model sizes, 10 workload scenarios, and up to
8 devices, classifying 732 individual operator executions.

\item \textbf{5-dimensional design space exploration} over core count,
HBM channels, L2 cache, systolic array dimensions, and SRAM, evaluating
5,082 feasible designs and identifying 79 Pareto-optimal configurations.

\item \textbf{Six actionable design guidelines} for next-generation LLM
accelerators, backed by quantitative analysis of parameter sensitivity,
area efficiency, and scaling behavior.
\end{enumerate}

%% ===========================================================================
\section{Background and Related Work}
\label{sec:background}

\subsection{LLM Inference Architecture}

A transformer-based LLM layer consists of multi-head attention (MHA)
and feed-forward network (FFN) blocks. The MHA block includes QKV
projection, $Q \cdot K^T$ attention scores, softmax, $A \cdot V$
context aggregation, and output projection. The FFN block includes
two large matrix multiplications with a GeLU activation between them.
Layer normalization is applied before each sub-block.

During prefill, the input sequence of length $s$ is processed in
parallel, producing matrices of shape $(b, s, d)$ where $b$ is the
batch size and $d$ is the model dimension. During decode, each forward
pass processes a single token, reducing the effective matrix dimensions
to $(b, 1, d)$ and shifting the bottleneck from compute to memory bandwidth.

\subsection{Roofline Performance Model}

The roofline model~\cite{williams2009roofline} characterizes operator
performance by comparing arithmetic intensity (AI = FLOPs/bytes) against
the machine's ridge point ($\pi/\beta$, where $\pi$ is peak compute and
$\beta$ is peak bandwidth). An operator with AI above the ridge point is
\emph{compute-bound}; otherwise it is \emph{memory-bound}. The attainable
performance is:

\begin{equation}
\text{Perf} = \min(\pi, \text{AI} \times \beta)
\end{equation}

And the execution time is:

\begin{equation}
T = \max\left(\frac{\text{FLOPs}}{\pi}, \frac{\text{Bytes}}{\beta}\right)
\end{equation}

\subsection{Related Work}

Prior work on GPU architecture for deep learning has focused on
training-oriented design~\cite{jouppi2017datacenter, jia2019dissecting},
while inference-specific studies have largely targeted CNN
workloads~\cite{chen2019eyeriss}. Recent attention to LLM inference
optimization has focused on algorithmic approaches (FlashAttention~\cite{dao2022flashattention},
PagedAttention~\cite{kwon2023efficient}, speculative
decoding~\cite{leviathan2023fast}) rather than architectural DSE.
LLMCompass~\cite{llmcompass} provides the first hardware-software
co-simulation framework specifically designed for LLM inference
architecture exploration, which we extend with systematic DSE methodology.

%% ===========================================================================
\section{Methodology}
\label{sec:methodology}

\subsection{Simulation Framework}

We use LLMCompass to model GPU systems at the operator level. The framework
models a \texttt{System} comprising one or more \texttt{Device}s connected
by an \texttt{InterConnect}. Each device contains a \texttt{ComputeModule}
(with configurable core count and systolic arrays), an \texttt{IOModule}
(with configurable L2 cache and HBM channels), and \texttt{Memory}.

For each operator, the framework computes:
\begin{itemize}
\item FLOP count and IO byte count based on tensor dimensions
\item Compute time = FLOPs / peak\_compute
\item Memory time = IO\_bytes / peak\_bandwidth
\item Latency = $\max(\text{compute\_time}, \text{memory\_time})$
\end{itemize}

A cost model estimates die area based on core count, systolic array size,
SRAM capacity, L2 cache size, HBM channels, and process node (7nm TSMC),
decomposed into compute chiplet area and IO die area.

\subsection{Workload Suite}

We profile 5 transformer models spanning 3 orders of magnitude in parameter
count (Table~\ref{tab:models}).

\begin{table}[t]
\centering
\caption{Model configurations used in profiling.}
\label{tab:models}
\begin{tabular}{lrrrr}
\toprule
\textbf{Model} & $d_{\text{model}}$ & $n_{\text{heads}}$ & $n_{\text{layers}}$ & \textbf{Params} \\
\midrule
GPT-2 Small  & 768    & 12  & 12  & 117M \\
GPT-2 XL     & 1600   & 25  & 48  & 1.5B \\
LLaMA-7B     & 4096   & 32  & 32  & 7B \\
LLaMA-70B    & 8192   & 64  & 80  & 70B \\
GPT-3 175B   & 12288  & 96  & 96  & 175B \\
\bottomrule
\end{tabular}
\end{table}

Each model is evaluated across 10 workload scenarios varying batch size
(1--64), sequence length (128--8192), phase (prefill/decode), and device
count (1/4/8), yielding 61 profiling tasks with 732 operator-level
measurements.

\subsection{Design Space Definition}

For DSE, we sweep 5 architectural parameters (Table~\ref{tab:dse_params})
across 3 device counts (1, 4, 8), generating 7,875 candidate designs.
Designs exceeding the 900~mm$^2$ reticle limit are pruned, leaving 5,082
feasible designs for evaluation.

\begin{table}[t]
\centering
\caption{DSE parameter space.}
\label{tab:dse_params}
\begin{tabular}{llr}
\toprule
\textbf{Parameter} & \textbf{Values} & \textbf{Count} \\
\midrule
Core count & 32, 48, 64, 96, 128, 192, 256 & 7 \\
HBM channels & 3, 4, 5, 6, 8 & 5 \\
L2 cache (MB) & 20, 40, 60, 80, 100 & 5 \\
Systolic array dim & 8, 16, 32 & 3 \\
SRAM per core (KB) & 32, 64, 128, 256, 512 & 5 \\
\midrule
Device counts & 1, 4, 8 & 3 \\
\bottomrule
\end{tabular}
\end{table}

% Baseline architecture: NVIDIA A100-like (GA100) with 108 SMs,
% 40MB L2, 5 HBM2e stacks, 32×32 systolic arrays.

\subsection{Evaluation Methodology}

Each design point is evaluated on 4 representative workloads:
LLaMA-7B and LLaMA-70B, each in prefill and decode phases, at batch
size 16 and sequence length 1024. We use three optimization objectives:
\begin{enumerate}
\item Minimize worst-case prefill latency (ms)
\item Minimize worst-case decode latency (ms)
\item Minimize die area (mm$^2$)
\end{enumerate}

Pareto-optimal designs are identified across all three objectives using
non-dominated sorting.

%% ===========================================================================
\section{Profiling Results: Bottleneck Analysis}
\label{sec:profiling}

\subsection{The Prefill-Decode Duality}

Table~\ref{tab:bottleneck} presents the operator-level bottleneck
classification across all 61 workload configurations. The most striking
finding is the \emph{complete asymmetry} between phases.

\begin{table}[t]
\centering
\caption{Operator bottleneck classification by phase.}
\label{tab:bottleneck}
\begin{tabular}{lrr}
\toprule
& \textbf{Prefill} & \textbf{Decode} \\
\midrule
Compute-bound operators & 114 (32.8\%) & 0 (0.0\%) \\
Memory-bound operators  & 234 (67.2\%) & 384 (100\%) \\
\midrule
Mean compute fraction   & 0.744 & 0.000 \\
\bottomrule
\end{tabular}
\end{table}

Decode is \textbf{universally memory-bound}: across all 384 operator
measurements spanning 5 models, batch sizes 1--64, sequence lengths
128--8192, and 1--8 devices, \emph{not a single operator} becomes
compute-bound during decode. This is not a matter of degree---the
compute fraction rounds to exactly zero.

\subsection{Operator-Level Invariants}

We observe that bottleneck classification is determined primarily by
\emph{operator type}, not workload parameters:

\begin{itemize}
\item \textbf{Always compute-bound in prefill}: FFN matmuls
(\texttt{h\_matmul1}, \texttt{h\_matmul2}) and QKV projection
(\texttt{qkv\_proj}, \texttt{h\_matmul0}). These operators have
arithmetic intensity $>$1500, far exceeding the ridge point ($\sim$180).

\item \textbf{Always memory-bound}: Attention operators
(\texttt{q\_mul\_k}, \texttt{a\_mul\_v}), vector operations
(\texttt{softmax}, \texttt{layernorm}, \texttt{gelu}), and
communication (\texttt{allreduce}). These operators remain
memory-bound even during prefill at batch size 64.
\end{itemize}

\subsection{Latency Dominance Analysis}

The FFN pathway dominates total latency:
\texttt{h\_matmul1} and \texttt{h\_matmul2} together account for
\textbf{48.2\%} of total inference time. The QKV projection adds
18.4\%, bringing the matmul-dominated fraction to 66.6\%. This
motivates prioritizing systolic array performance over vector unit
or memory subsystem optimization for prefill workloads.

%% ===========================================================================
\section{Design Space Exploration Results}
\label{sec:dse}

\subsection{Parameter Sensitivity}

Table~\ref{tab:sensitivity} presents the key finding of our sensitivity
analysis. We measure the ``range percent''---the percentage difference
between the best and worst mean objective value as each parameter is swept
while others are averaged.

\begin{table}[t]
\centering
\caption{Parameter sensitivity ranking by objective.}
\label{tab:sensitivity}
\begin{tabular}{lrrr}
\toprule
\textbf{Parameter} & \textbf{Prefill (\%)} & \textbf{Decode (\%)} & \textbf{Area (\%)} \\
\midrule
Systolic array dim & \textbf{1074.8} & 0.2 & 4.9 \\
Core count         & 206.2 & 17.5 & \textbf{112.9} \\
HBM channels       & 3.1 & \textbf{35.2} & 10.2 \\
L2 cache (MB)      & 17.4 & 0.3 & 31.7 \\
SRAM (KB)          & 8.1 & 0.2 & 5.0 \\
\bottomrule
\end{tabular}
\end{table}

Three critical observations emerge:

\paragraph{Systolic array dimension dominates prefill.}
The systolic array dimension exhibits 1074.8\% range on prefill latency---
5.2$\times$ more impactful than core count (206.2\%)---while consuming
only 4.9\% of area variation. This occurs because systolic array throughput
scales \emph{quadratically} with dimension ($O(n^2)$), while core count
scales linearly. Going from 8$\times$8 to 16$\times$16 systolic arrays
reduces mean prefill latency by 3.75$\times$ at a cost of only 7~mm$^2$,
yielding an area efficiency of \textbf{2,587~ms/mm$^2$}---31.7$\times$ more
efficient than adding cores (81.6~ms/mm$^2$).

\paragraph{HBM bandwidth is the sole decode lever.}
HBM channel count accounts for 35.2\% of decode variation. All other
parameters \emph{combined} contribute $<$1\%. Systolic array dimension---
the most impactful prefill parameter---has only 0.16\% decode impact.
This means prefill and decode optimization are \emph{almost completely
orthogonal}.

\paragraph{On-chip memory is counter-productive.}
L2 cache shows 17.4\% prefill impact, but in the \emph{wrong direction}:
larger L2 increases latency, likely due to higher data movement costs
in the roofline model. Meanwhile, its decode impact is 0.3\%
(negligible), yet it consumes 31.7\% of area variation.

\subsection{Pareto Frontier Analysis}

Of 5,082 feasible designs, 79 are Pareto-optimal across prefill latency,
decode latency, and die area. The frontier reveals a clear structure
(Figure~\ref{fig:pareto}):

\begin{itemize}
\item \textbf{4-device designs} (48/79 Pareto points) achieve the
best decode latency (34.5~ms) but higher prefill (766.7~ms minimum).
\item \textbf{8-device designs} (31/79 Pareto points) achieve the
best prefill latency (550.5~ms) but worse decode (47.7~ms minimum).
\item \textbf{1-device designs}: none appear on the Pareto frontier.
\end{itemize}

The Pearson correlation between prefill and decode latency on the Pareto
frontier is $r = -0.294$, confirming that these objectives are
structurally anti-correlated.

\begin{figure}[t]
\centering
\includegraphics[width=0.95\columnwidth]{figures/pareto_placeholder}
\caption{Pareto frontier: prefill vs decode latency, colored by die area.
Pareto-optimal designs (large markers) span from compute-optimized (left)
to decode-optimized (right) configurations. See supplementary for full plots.}
\label{fig:pareto}
\end{figure}

\subsection{Area-Performance Tradeoffs}

Table~\ref{tab:area_eff} shows the marginal return per mm$^2$ for each
parameter increment.

\begin{table}[t]
\centering
\caption{Marginal prefill latency reduction per mm$^2$ of die area.}
\label{tab:area_eff}
\begin{tabular}{lrrr}
\toprule
\textbf{Parameter Step} & $\Delta$\textbf{Lat (ms)} & $\Delta$\textbf{Area} & \textbf{ms/mm$^2$} \\
\midrule
SA dim 8$\to$16     & $-$18,071 & +7.0    & \textbf{2,587} \\
SA dim 16$\to$32    & $-$4,477  & +21.1   & 212 \\
Cores 32$\to$48     & $-$6,283  & +77.0   & 82 \\
Cores 48$\to$64     & $-$3,130  & +77.0   & 41 \\
Cores 96$\to$128    & $-$645    & +85.6   & 8 \\
Cores 128$\to$192   & \textbf{+1,397} & +94.7 & \textbf{$-$15} \\
\bottomrule
\end{tabular}
\end{table}

The most significant finding is that systolic array scaling is
\textbf{31.7$\times$ more area-efficient} than core scaling for prefill
optimization. Furthermore, core counts beyond 128 yield \emph{negative
returns}---latency increases by 1,397~ms while consuming 94.7~mm$^2$
additional area.

%% ===========================================================================
\section{Multi-Device Scaling Analysis}
\label{sec:scaling}

\subsection{Prefill Scaling: Near-Linear}

Prefill exhibits good scaling behavior with multi-device tensor
parallelism: 86.9\% efficiency at 4 devices and 78.2\% at 8 devices.
This is expected because prefill operators have high arithmetic intensity
and large tensor dimensions, allowing effective parallelization.

\subsection{Decode Scaling: The Scaling Wall}

Decode scaling is severely sub-linear and can be counter-productive:

\begin{itemize}
\item 1$\to$4 devices: 38.8\% efficiency
\item 1$\to$8 devices: \textbf{15.0\%} efficiency
\end{itemize}

In concrete terms, for LLaMA-70B decode (bs=16, sl=1024): going from
4 to 8 devices \emph{increases} latency from 0.51~ms to 0.64~ms.
The culprit is communication overhead: at 8 devices, allreduce consumes
37--40\% of total decode time.

This has profound implications for inference serving: \textbf{4-device
clusters are strictly superior to 8-device clusters for decode-heavy
autoregressive generation.}

%% ===========================================================================
\section{Design Guidelines}
\label{sec:guidelines}

Synthesizing our findings, we propose six design guidelines for
next-generation LLM inference accelerators:

\begin{enumerate}
\item \textbf{Prioritize systolic array dimensions over core count.}
A 32$\times$32 systolic array is 31.7$\times$ more area-efficient than
adding cores. Current GPU architectures significantly underinvest in
per-core compute density.

\item \textbf{Maximize HBM bandwidth per device.}
HBM channel count is the \emph{only} lever for decode acceleration
(35.2\% impact vs $<$1\% for all other parameters). Every additional
HBM channel reduces decode latency by 2--6~ms at a cost of $\sim$11~mm$^2$.

\item \textbf{Cap core count at 96--128.}
Beyond 128 cores, parallelization overhead exceeds compute gains,
resulting in negative performance returns while consuming 12\%+
additional area.

\item \textbf{Optimize for 4-device, not 8-device decode serving.}
4-device configurations achieve 1.38$\times$ better decode latency
than 8-device due to reduced communication overhead.

\item \textbf{Minimize L2 cache and SRAM for inference.}
On-chip memory consumes 31.7\% of area variation with negligible-to-negative
performance impact for LLM inference. Area savings should be reallocated
to either HBM channels or systolic array capacity.

\item \textbf{Consider prefill-decode architectural disaggregation.}
The anti-correlation ($r = -0.29$) between prefill and decode
optimization on the Pareto frontier suggests that no single architecture
can optimally serve both. Disaggregated prefill-decode serving
(separate chiplets or server pools) can eliminate this tradeoff.
\end{enumerate}

%% ===========================================================================
\section{Discussion}
\label{sec:discussion}

\paragraph{Validity of the roofline model.}
Our analysis relies on the roofline model, which assumes perfect pipelining
of compute and memory operations. Real GPU execution involves additional
overheads (warp scheduling, bank conflicts, cache misses) that are not captured.
However, the roofline model provides accurate \emph{relative rankings} of
architectural configurations, which is sufficient for DSE purposes.

\paragraph{The 128-core cliff.}
The performance degradation beyond 128 cores may partly reflect the cost
model's assumption about parallelism overhead rather than a fundamental
architectural limit. Real-world implementations with more sophisticated
work distribution may push this limit higher, but the \emph{diminishing
returns} pattern is robust.

\paragraph{Missing workload dimensions.}
Our study covers standard dense transformer architectures. Mixture-of-Experts
(MoE) models, which activate only a subset of parameters per token, may shift
the compute-memory balance. Similarly, very long sequences ($>$8K tokens)
with KV-cache management may change the on-chip memory calculus.

\paragraph{Energy and power.}
Our analysis focuses on area-performance tradeoffs. Energy-performance
tradeoffs, which are increasingly critical for datacenter deployments,
may yield different conclusions---particularly for the on-chip memory
parameters that consume area but might reduce data movement energy.

%% ===========================================================================
\section{Conclusion}
\label{sec:conclusion}

Through systematic profiling and design space exploration of GPU
architectures for LLM inference, we demonstrate that the prefill and
decode phases present fundamentally opposed architectural requirements.
Our analysis reveals that systolic array dimensions are the most
underpriced resource in current GPU designs (31.7$\times$ more
area-efficient than cores), while on-chip memory is the most overpriced.
HBM bandwidth is the sole determinant of decode performance, and
multi-device scaling hits a wall at 4 devices for decode workloads.
These findings motivate a paradigm shift toward heterogeneous
prefill-decode architectures where compute-dense prefill engines are
paired with bandwidth-optimized decode engines, potentially
disaggregated across different chiplets or server pools.

%% ===========================================================================
\bibliographystyle{ACM-Reference-Format}

\begin{thebibliography}{10}

\bibitem{brown2020language}
T.~Brown et~al.
\newblock Language models are few-shot learners.
\newblock In \emph{NeurIPS}, 2020.

\bibitem{touvron2023llama}
H.~Touvron et~al.
\newblock LLaMA: Open and efficient foundation language models.
\newblock \emph{arXiv:2302.13971}, 2023.

\bibitem{williams2009roofline}
S.~Williams, A.~Waterman, and D.~Patterson.
\newblock Roofline: An insightful visual performance model for multicore
  architectures.
\newblock \emph{CACM}, 52(4):65--76, 2009.

\bibitem{llmcompass}
H.~Zhang et~al.
\newblock LLMCompass: Enabling efficient hardware design for large language
  model inference.
\newblock In \emph{ISCA}, 2024.

\bibitem{jouppi2017datacenter}
N.~Jouppi et~al.
\newblock In-datacenter performance analysis of a tensor processing unit.
\newblock In \emph{ISCA}, 2017.

\bibitem{jia2019dissecting}
Z.~Jia, M.~Maggioni, B.~Staiger, and D.~Scarpazza.
\newblock Dissecting the NVIDIA Volta GPU architecture via microbenchmarking.
\newblock \emph{arXiv:1804.06826}, 2019.

\bibitem{chen2019eyeriss}
Y.-H.~Chen et~al.
\newblock Eyeriss v2: A flexible accelerator for emerging deep neural networks.
\newblock \emph{JETCAS}, 2019.

\bibitem{dao2022flashattention}
T.~Dao et~al.
\newblock FlashAttention: Fast and memory-efficient exact attention with
  IO-awareness.
\newblock In \emph{NeurIPS}, 2022.

\bibitem{kwon2023efficient}
W.~Kwon et~al.
\newblock Efficient memory management for large language model serving with
  PagedAttention.
\newblock In \emph{SOSP}, 2023.

\bibitem{leviathan2023fast}
Y.~Leviathan, M.~Kalman, and Y.~Matias.
\newblock Fast inference from transformers via speculative decoding.
\newblock In \emph{ICML}, 2023.

\end{thebibliography}

\end{document}
